\chapter{Vetores, retas e planos no espaço}\label{ch:g12:space}

A geometria tridimensional torna-se calculável assim que dispomos de
\omterm{def:g11:vect:vector}{vetores} e \omterm{def:g10:coordgeom:system}{coordenadas}: as retas ganham representações paramétricas, os
planos ganham \omterm{def:g10:algebra:equation}{equações} cartesianas, e o \omterm{def:g12:space:dot}{produto escalar} mede ângulos e
distâncias. Este capítulo constrói essa caixa de ferramentas e a usa para
resolver os problemas clássicos de \omterm{def:g10:numbers:interunion}{interseção} e de distância.

\section{Vetores no espaço}

Os \omterm{def:g11:vect:vector}{vetores} no espaço obedecem às mesmas regras que no plano: somam-se,
multiplicam-se por escalares, e $\vect{AB} = \vect{CD}$ exatamente quando
$ABDC$ é um paralelogramo. Em um \omterm{def:g10:coordgeom:system}{sistema de coordenadas}
$(O; \vec\imath, \vec\jmath, \vec k)$, todo \omterm{def:g11:vect:vector}{vetor} tem \omterm{def:g10:coordgeom:system}{coordenadas}
$(x, y, z)$.

\begin{definition}[Colinearidade, coplanaridade]\label{def:g12:space:collinear}
Dois \omterm{def:g11:vect:vector}{vetores} são \emph{\omterm{def:g11:vect:collinear}{colineares}} se um é múltiplo do outro. Três \omterm{def:g11:vect:vector}{vetores}
$\vec u, \vec v, \vec w$ são \emph{coplanares}\index{coplanares} se um
deles pode ser escrito como combinação dos outros dois, digamos
$\vec w = a\vec u + b\vec v$.
\end{definition}

\begin{proposition}\label{prop:g12:space:basis}
Se $\vec u, \vec v, \vec w$ não são \omterm{def:g12:space:collinear}{coplanares}, todo \omterm{def:g11:vect:vector}{vetor} do espaço se
decompõe de maneira única como $x\vec u + y\vec v + z\vec w$.
\end{proposition}

\begin{proof}[Ideia da demonstração]
Pela ponta do \omterm{def:g11:vect:vector}{vetor} a decompor, trace a reta paralela a $\vec w$; ela
encontra o plano de $\vec u, \vec v$ em um único ponto, que separa o \omterm{def:g11:vect:vector}{vetor}
em uma componente nesse plano (de maneira única $x\vec u + y\vec v$,
geometria plana) e uma componente na direção de $\vec w$. Unicidade: a
diferença de duas decomposições exprimiria um dos \omterm{def:g11:vect:vector}{vetores} em \omterm{def:g11:func:function}{função} dos
outros dois, contradizendo a não coplanaridade.
\end{proof}

\section{Retas e planos}

\begin{definition}[Representação paramétrica de uma reta]\label{def:g12:space:line}
A reta que passa por $A(x_A, y_A, z_A)$ com \omterm{def:g11:vect:direction}{vetor diretor}
$\vec u (a, b, c) \neq \vec 0$ é o conjunto dos pontos
\[
M(x_A + ta,\; y_A + tb,\; z_A + tc), \qquad t \in \R .
\]
\end{definition}

\begin{definition}[Plano]\label{def:g12:space:plane}
O plano que passa por $A$ e é dirigido por dois \omterm{def:g11:vect:vector}{vetores} não \omterm{def:g12:space:collinear}{colineares}
$\vec u, \vec v$ é o conjunto dos pontos $M$ com
$\vect{AM} = s\vec u + t\vec v$, $(s, t) \in \R^2$.
\end{definition}

\begin{omfigure}
\begin{tikzpicture}[scale=1.05]
  \draw (-0.5,-0.25) -- (3.4,1.7);
  \draw[-Stealth, omDef, thick] (1,0.5) -- (1.9,0.95)
    node[midway, above left, font=\small] {$\vec u$};
  \fill (0.1,0.05) circle (1.2pt) node[below right, font=\small] {$t=-1$};
  \fill (1,0.5) circle (1.2pt) node[above left, font=\small] {$A$};
  \fill (1.9,0.95) circle (1.2pt) node[below right, font=\small] {$t=1$};
  \fill (2.8,1.4) circle (1.2pt) node[below right, font=\small] {$t=2$};
\end{tikzpicture}
\qquad\quad
\begin{tikzpicture}[scale=1.05]
  \draw (0,0) -- (4,0.5) -- (5.2,1.9) -- (1.2,1.4) -- cycle;
  \node[font=\small] at (4.75,1.6) {$\mathcal P$};
  \coordinate (A) at (1.3,0.55);
  \coordinate (M) at (3.77,1.38);
  \draw[-Stealth, omDef, thick] (A) -- ++(1.31,0.16)
    node[midway, below, font=\small] {$\vec u$};
  \draw[-Stealth, omDef, thick] (A) -- ++(0.42,0.49)
    node[midway, left, font=\small] {$\vec v$};
  \draw[black, dashed, thin] (3.27,0.79) -- (M) -- (1.80,1.14);
  \fill (A) circle (1.2pt) node[below left, font=\small] {$A$};
  \fill (M) circle (1.2pt) node[above, font=\small] {$M$};
\end{tikzpicture}

{\small À esquerda: a reta que passa por $A$ com direção $\vec u$,
graduada pelo parâmetro $t$. À direita: o plano que passa por $A$ dirigido
por $\vec u$ e $\vec v$; todo ponto $M$ do plano é atingido como
$\vect{AM} = s\vec u + t\vec v$.}
\end{omfigure}

\begin{proposition}[Posições relativas]\label{prop:g12:space:positions}
Dois planos distintos ou são paralelos, ou se cortam segundo uma reta. Uma
reta e um plano ou são paralelos (podendo a reta estar contida no plano),
ou se cortam em um único ponto. Duas retas no espaço podem ser
concorrentes, estritamente paralelas, coincidentes ou \emph{reversas} (não
\omterm{def:g12:space:collinear}{coplanares}).
\end{proposition}

\begin{proof}[Esboço]
São afirmações de incidência; cada análise de casos se reduz, em
\omterm{def:g10:coordgeom:system}{coordenadas}, a resolver um \omterm{def:g10:lines:system}{sistema linear} e contar suas soluções --- veja
o \cref{met:g12:space:intersections}.
\end{proof}

\section{Produto escalar}

\begin{definition}[Produto escalar]\label{def:g12:space:dot}
O \emph{produto escalar}\index{produto escalar} de $\vec u$ e $\vec v$ é
\[
\vec u \cdot \vec v = \tfrac12\left(\norm{\vec u + \vec v}^2
- \norm{\vec u}^2 - \norm{\vec v}^2\right),
\]
em que $\norm{\vec u}$ é o comprimento de $\vec u$. Se os dois \omterm{def:g11:vect:vector}{vetores} são
não nulos, $\vec u \cdot \vec v = \norm{\vec u}\,\norm{\vec v}\cos\theta$,
sendo $\theta$ o ângulo entre eles; e, em um \omterm{def:g10:coordgeom:system}{sistema de coordenadas}
\emph{\omterm{def:g10:coordgeom:system}{ortonormal}},
\[
\vec u \cdot \vec v = xx' + yy' + zz' .
\]
Dois \omterm{def:g11:vect:vector}{vetores} são \emph{\omterm{def:g11:scal:orthogonal}{ortogonais}} se seu produto escalar é $0$.
\end{definition}

\begin{proposition}[Propriedades]\label{prop:g12:space:dotprops}
O \omterm{def:g12:space:dot}{produto escalar} é simétrico ($\vec u\cdot\vec v = \vec v\cdot\vec u$),
bilinear ($(a\vec u + b\vec v)\cdot \vec w = a\,\vec u\cdot\vec w +
b\,\vec v\cdot \vec w$), e
$\vec u \cdot \vec u = \norm{\vec u}^2 \geq 0$.
\end{proposition}

\begin{proof}
Em um \omterm{def:g10:coordgeom:system}{sistema ortonormal}, as três propriedades são imediatas na fórmula
$xx' + yy' + zz'$, que por sua vez decorre da fórmula de definição e do
cálculo pitagórico dos comprimentos:
$\norm{\vec u}^2 = x^2 + y^2 + z^2$.
\end{proof}

\begin{definition}[Vetor normal]\label{def:g12:space:normal}
Um \emph{vetor normal}\index{vetor normal} de um plano $\mathcal P$ é um
\omterm{def:g11:vect:vector}{vetor} não nulo \omterm{def:g11:scal:orthogonal}{ortogonal} a todo \omterm{def:g11:vect:vector}{vetor} contido em $\mathcal P$ (basta que
seja \omterm{def:g11:scal:orthogonal}{ortogonal} a duas direções não \omterm{def:g12:space:collinear}{colineares} de $\mathcal P$).
\end{definition}

\begin{omfigure}
\begin{tikzpicture}[scale=1.1]
  \draw (0,0) -- (4,0.5) -- (5.2,1.9) -- (1.2,1.4) -- cycle;
  \node[font=\small] at (4.75,1.6) {$\mathcal P$};
  \coordinate (H) at (2.6,0.95);
  \coordinate (M0) at (2.86,3.03);
  \draw[-Stealth, omThm, thick] (H) -- (2.73,1.99)
    node[midway, right=1pt, font=\small] {$\vec n$};
  \draw[black, dashed, thin] (2.73,1.99) -- (M0);
  \draw[thin] (2.79,0.974) -- (2.815,1.174) -- (2.625,1.15);
  \draw[black, dashed, thin] (H) -- (1.1,0.55);
  \fill (H) circle (1.2pt) node[below right, font=\small] {$H$};
  \fill (M0) circle (1.2pt) node[above, font=\small] {$M_0$};
  \fill (1.1,0.55) circle (1.2pt) node[below left, font=\small] {$A$};
\end{tikzpicture}

{\small Um \omterm{def:g12:space:normal}{vetor normal} $\vec n$ do plano $\mathcal P$ é \omterm{def:g11:scal:orthogonal}{ortogonal} a toda
direção de $\mathcal P$. O ponto $H$, pé da perpendicular baixada de
$M_0$, realiza a distância de $M_0$ ao plano.}
\end{omfigure}

\begin{theorem}[Equação cartesiana de um plano]\label{thm:g12:space:planeeq}
Em um \omterm{def:g10:coordgeom:system}{sistema de coordenadas} \omterm{def:g10:coordgeom:system}{ortonormal}, o plano que passa por $A$ com
\omterm{def:g12:space:normal}{vetor normal} $\vec n(a, b, c)$ tem \omterm{def:g10:algebra:equation}{equação}
\[
a x + b y + c z + d = 0,
\]
em que $d = -(ax_A + by_A + cz_A)$. Reciprocamente, toda \omterm{def:g10:algebra:equation}{equação} desse
tipo com $(a,b,c) \neq (0,0,0)$ define um plano de \omterm{def:g12:space:normal}{vetor normal}
$(a, b, c)$.
\end{theorem}

\begin{proof}
$M \in \mathcal P \iff \vect{AM} \perp \vec n \iff
a(x - x_A) + b(y - y_A) + c(z - z_A) = 0$, o que, desenvolvido, dá a
\omterm{def:g10:algebra:equation}{equação}. Reciprocamente, tome qualquer ponto $A$ que satisfaça a \omterm{def:g10:algebra:equation}{equação};
o mesmo cálculo, ao contrário, mostra que o conjunto solução é
$\{M : \vect{AM}\cdot\vec n = 0\}$, o plano que passa por $A$ e é normal a
$\vec n$.
\end{proof}

\begin{method}[Interseções na prática]\label{met:g12:space:intersections}
\begin{itemize}
  \item \emph{Reta $\cap$ plano}: substitua as \omterm{def:g10:coordgeom:system}{coordenadas} paramétricas da
        reta na \omterm{def:g10:algebra:equation}{equação} do plano; resolva em $t$ (uma solução: um ponto;
        $0 = $ não nulo: paralelismo; $0 = 0$: reta contida).
  \item \emph{Plano $\cap$ plano}: resolva o sistema das duas \omterm{def:g10:algebra:equation}{equações},
        parametrizando por uma coordenada livre; a solução é uma reta (ou
        os planos são paralelos, quando os \omterm{def:g12:space:normal}{vetores normais} são
        \omterm{def:g12:space:collinear}{colineares}).
  \item \emph{Testes de ortogonalidade}: reta $\perp$ plano se, e somente
        se, sua direção é \omterm{def:g11:vect:collinear}{colinear} com a normal; dois planos são
        perpendiculares se, e somente se, suas normais são \omterm{def:g11:scal:orthogonal}{ortogonais}.
\end{itemize}
\end{method}

\begin{proposition}[Distância de um ponto a um plano]\label{prop:g12:space:distance}
\index{distância!de um ponto a um plano}
Em um \omterm{def:g10:coordgeom:system}{sistema ortonormal}, a distância de $M_0(x_0, y_0, z_0)$ ao plano
$\mathcal P : ax + by + cz + d = 0$ é
\[
\operatorname{dist}(M_0, \mathcal P)
= \frac{\abs{ax_0 + by_0 + cz_0 + d}}{\sqrt{a^2 + b^2 + c^2}}.
\]
\end{proposition}

\begin{proof}
Seja $H$ a projeção \omterm{def:g11:scal:orthogonal}{ortogonal} de $M_0$ sobre $\mathcal P$: o pé da
perpendicular, de modo que $\vect{HM_0}$ é \omterm{def:g11:vect:collinear}{colinear} com a normal unitária
$\frac{\vec n}{\norm{\vec n}}$, e a distância é
$\abs{\vect{HM_0}\cdot \frac{\vec n}{\norm{\vec n}}}$. Para qualquer
$H \in \mathcal P$,
$\vect{HM_0}\cdot\vec n = ax_0 + by_0 + cz_0 - (ax_H + by_H + cz_H)
= ax_0 + by_0 + cz_0 + d$ (usando a \omterm{def:g10:algebra:equation}{equação} em $H$), donde a fórmula,
depois de dividir por $\norm{\vec n} = \sqrt{a^2+b^2+c^2}$.
\end{proof}

\begin{example}\label{ex:g12:space:distance}
Distância da origem ao plano $x + 2y + 2z - 6 = 0$:
$\frac{\abs{-6}}{\sqrt{1 + 4 + 4}} = \frac63 = 2$.
\end{example}

\section{Exercícios}

Em todos os exercícios, o \omterm{def:g10:coordgeom:system}{sistema de coordenadas} é \omterm{def:g10:coordgeom:system}{ortonormal}.

\begin{exercise}[$\star$]\label{exo:g12:space:1}
Sejam $A(1, 0, 2)$, $B(3, 1, 1)$ e $C(2, -1, 3)$. Os pontos $A$, $B$, $C$
estão alinhados? Dê uma representação paramétrica da reta $(AB)$.
\end{exercise}

\begin{exercise}[$\star$]\label{exo:g12:space:2}
Determine uma \omterm{def:g10:algebra:equation}{equação} cartesiana do plano que passa por $A(1, 1, 0)$ com
\omterm{def:g12:space:normal}{vetor normal} $\vec n(2, -1, 3)$. O ponto $B(0, 2, 1)$ pertence a ele?
\end{exercise}

\begin{exercise}[$\star$]\label{exo:g12:space:3}
Calcule o ângulo em $A$ (ao grau mais próximo) do triângulo de vértices
$A(0,0,0)$, $B(1, 1, 0)$ e $C(1, 0, 1)$.
\end{exercise}

\begin{exercise}[$\star\star$]\label{exo:g12:space:4}
Considere a reta $\mathcal D$ que passa por $A(1, 2, 0)$ com direção
$\vec u(1, -1, 2)$, e o plano $\mathcal P : 2x + y - z + 1 = 0$.
Determine $\mathcal D \cap \mathcal P$.
\end{exercise}

\begin{exercise}[$\star\star$]\label{exo:g12:space:5}
Sejam $\mathcal P : x - y + 2z = 1$ e $\mathcal Q : 2x + y + z = 4$.
Mostre que $\mathcal P$ e $\mathcal Q$ se cortam segundo uma reta e dê uma
representação paramétrica dela.
\end{exercise}

\begin{exercise}[$\star\star$]\label{exo:g12:space:6}
Mostre que as retas
\[
\mathcal D_1 : (x, y, z) = (1 + t,\; 2t,\; 3 - t), \qquad
\mathcal D_2 : (x, y, z) = (2 + s,\; 1 + s,\; 1 + 2s)
\]
são reversas (não \omterm{def:g12:space:collinear}{coplanares}).
\end{exercise}

\begin{exercise}[$\star\star$]\label{exo:g12:space:7}
Seja $\mathcal P : 2x - y + 2z - 3 = 0$.
\begin{enumerate}
  \item Calcule a distância de $M_0(3, 1, 1)$ a $\mathcal P$.
  \item Determine as \omterm{def:g10:coordgeom:system}{coordenadas} da projeção \omterm{def:g11:scal:orthogonal}{ortogonal} $H$ de $M_0$ sobre
        $\mathcal P$ e confira a distância $M_0H$.
\end{enumerate}
\end{exercise}

\begin{exercise}[$\star\star\star$]\label{exo:g12:space:8}
O cubo $ABCDEFGH$ tem aresta $1$; escolha \omterm{def:g10:coordgeom:system}{coordenadas} de modo que
$A(0,0,0)$, $B(1,0,0)$, $D(0,1,0)$, $E(0,0,1)$ (com
$C = B + D$, $F = B + E$, $G = B + D + E$, $H = D + E$ como \omterm{def:g11:vect:vector}{vetores} a
partir de $A$).
\begin{enumerate}
  \item Mostre que a diagonal $(AG)$ é \omterm{def:g11:scal:orthogonal}{ortogonal} ao plano $(BDE)$.
  \item Calcule a distância de $A$ ao plano $(BDE)$ e o ponto em que
        $(AG)$ o atravessa.
\end{enumerate}
\end{exercise}

\begin{exercise}[$\star\star\star$]\label{exo:g12:space:9}
\emph{(Volume de um tetraedro.)} Sejam $A(1,1,1)$, $B(2,1,0)$, $C(0,2,1)$
e $D(2,2,2)$.
\begin{enumerate}
  \item Verifique que $\vect{AB}$, $\vect{AC}$, $\vect{AD}$ não são
        \omterm{def:g12:space:collinear}{coplanares} (os pontos formam um tetraedro de verdade).
  \item Encontre um \omterm{def:g11:vect:vector}{vetor} $\vec n$ \omterm{def:g11:scal:orthogonal}{ortogonal} a $\vect{BC}$ e a
        $\vect{BD}$, e deduza uma \omterm{def:g10:algebra:equation}{equação} cartesiana do plano $(BCD)$.
  \item Calcule a distância de $A$ ao plano $(BCD)$ e a área do triângulo
        $BCD$, usando a fórmula
        $\text{Área} = \frac12\sqrt{\norm{\vec u}^2\norm{\vec v}^2 -
        (\vec u \cdot \vec v)^2}$ para o triângulo gerado por
        $\vec u, \vec v$.
  \item Deduza o volume do tetraedro
        ($V = \frac13 \times \text{área da base} \times \text{altura}$).
\end{enumerate}
\end{exercise}

\section{Problema: fatiando o cubo}

\begin{problem}[{Problema de fim de semana --- um plano corta um
cubo em um hexágono perfeito, um tetraedro regular se esconde nos
cantos, e o ângulo da diagonal dá forma a todo diamante}]
\label{pb:g12:space:1}
Fatie um cubo e você espera quadrados e retângulos --- e, no
entanto, um corte famoso produz um \emph{hexágono regular}, e
quatro cortes espertos nos cantos deixam para trás um
\emph{tetraedro regular}. As duas surpresas, e o ângulo de
$109.5^\circ$ que os átomos de carbono reverenciam, caem diante
da caixa de ferramentas deste capítulo: retas paramétricas,
\omterm{def:g10:algebra:equation}{equações} de planos e o \omterm{def:g12:space:dot}{produto escalar}
(\cref{thm:g12:space:planeeq},
\cref{prop:g12:space:distance}). Use por toda parte o cubo do
\cref{exo:g12:space:8}: aresta $1$, $A(0,0,0)$, $B(1,0,0)$,
$D(0,1,0)$, $E(0,0,1)$, $C(1,1,0)$, $F(1,0,1)$, $H(0,1,1)$,
$G(1,1,1)$.

\medskip
\noindent\textbf{Parte I --- Fluência.}
\begin{enumerate}
  \item Escreva a \omterm{def:g10:algebra:equation}{equação} do plano que passa por $B(1,0,0)$ com
        \omterm{def:g12:space:normal}{vetor normal} $(1,1,1)$, uma representação paramétrica da
        reta que passa por $A$ dirigida por $(1,1,1)$, e o ponto
        de \omterm{def:g10:numbers:interunion}{interseção} dos dois.
  \item Calcule a distância de $A$ a esse plano
        (\cref{prop:g12:space:distance}).
  \item Os \omterm{def:g11:vect:vector}{vetores} $(1,-1,0)$ e $(1,1,-2)$ são \omterm{def:g11:scal:orthogonal}{ortogonais}?
  \item Posição dos planos $x + y + z = 1$ e
        $2x + 2y + 2z = 5$? E da reta que passa por $A$ dirigida
        por $(1,1,0)$ em relação ao primeiro plano?
  \item Liste os \omterm{prop:g10:coordgeom:midpoint}{pontos médios} das seis arestas
        $[BC], [CD], [DH], [HE], [EF], [FB]$ do cubo e
        verifique que todos os seis satisfazem a \omterm{def:g10:algebra:equation}{equação}
        $x + y + z = \frac32$.
\end{enumerate}

\medskip
\noindent\textbf{Parte II --- A seção hexagonal.}
Seja $P$ o plano $x + y + z = \frac32$.
\begin{enumerate}[resume]
  \item Mostre que $P$ é perpendicular à grande diagonal
        $(AG)$ e passa pelo centro do cubo.
  \item Calcule as distâncias entre \omterm{prop:g10:coordgeom:midpoint}{pontos médios} consecutivos
        da questão 5 ao longo da seção: o que você encontra?
  \item Calcule o ângulo do polígono da seção em um dos
        vértices (dois vetores-aresta consecutivos e um \omterm{def:g12:space:dot}{produto
        escalar}). Conclua: a seção é um \emph{hexágono
        regular}.
  \item Calcule o perímetro e a área do hexágono, e compare a
        área com a de uma face do cubo. (Para o registro: a
        maior seção plana do cubo unitário é o retângulo
        diagonal $1 \times \sqrt2$, de área $\sqrt2$ --- o
        hexágono regular fica com a prata.)
  \item Onde a diagonal $(AG)$ atravessa $P$? Verifique que o
        ponto de travessia é o \omterm{prop:g10:coordgeom:midpoint}{ponto médio} de $[AG]$.
  \item Deslize o corte: descreva as seções
        $x + y + z = c$ à medida que $c$ cresce de $0$ a $3$ ---
        calcule a seção para $c = \frac12$ (qual polígono, de
        que tamanho?) e narre a metamorfose que passa pelo seu
        hexágono em $c = \frac32$.
  \item Por que nenhum plano pode cortar um cubo segundo um
        polígono de \emph{sete} lados? (Onde tem de estar cada
        lado da seção?)
\end{enumerate}

\medskip
\noindent\textbf{Parte III --- O tetraedro escondido.}
\begin{enumerate}[resume]
  \item Calcule o volume do tetraedro de canto $ABDE$ (base,
        altura e a fórmula do um terço que o volume do ensino
        fundamental admitiu --- o \cref{exo:g12:space:9}
        calcula volumes assim em geral).
  \item Os quatro vértices $B$, $D$, $E$, $G$: calcule as seis
        distâncias duas a duas e conclua que $BDEG$ é um
        tetraedro \emph{regular}. Depois obtenha seu volume
        subtraindo do cubo os tetraedros de canto.
  \item Calcule a distância do centro do cubo ao plano $(BDE)$.
  \item O ângulo do diamante: calcule o ângulo entre as
        diagonais $\vect{AG}$ e $\vect{BH}$ do cubo. Seu
        suplemento, cerca de $109.5^\circ$, é o ângulo entre as
        ligações no diamante e no metano --- por que quatro
        pares de elétrons em torno de um átomo de carbono
        escolhem os cantos do tetraedro da questão 14?
\end{enumerate}

\medskip
\noindent\textbf{Parte IV --- Genuinamente tridimensional.}
\begin{enumerate}[resume]
  \item Encontre onde a reta que passa por $B(1,0,0)$ dirigida
        por $(0,1,1)$ atravessa o plano $P$ do hexágono.
  \item Mostre que as retas $(AB)$ e $(EG)$ não se encontram
        nem são paralelas: retas \emph{reversas}, o fenômeno
        que não pode ocorrer em um plano.
  \item Calcule a projeção \omterm{def:g11:scal:orthogonal}{ortogonal} de $B$ sobre $P$ (caminhe
        de $B$ ao longo da normal até o plano) e verifique que
        seu ponto satisfaz a \omterm{def:g10:algebra:equation}{equação} de $P$.
  \item Finale --- o kit do agrimensor do espaço: retas
        paramétricas para atravessar, \omterm{def:g12:space:normal}{vetores normais} para
        ângulos e distâncias, \omterm{def:g10:algebra:equation}{equações} de planos para seções; e
        o cubo como o pátio de recreio em que produziram um
        hexágono, um tetraedro regular, o ângulo do diamante e
        um par de retas reversas. Uma frase para cada.
\end{enumerate}
\end{problem}
